% !TeX root = ../main.tex
\section{Introduction}
\label{sec:intro}

Application-layer encryption and modern transport protocols have reduced the
effectiveness of deep packet inspection (DPI) for identifying user activities.
QUIC, now widely deployed for web and video services, encrypts most protocol
headers and all payloads. Nevertheless, \emph{flow metadata}---sequence of packet
sizes, directions, and inter-packet times (IPTs)---is still visible on backbone and
access links. Prior work shows that machine learning models can infer application
classes from such metadata with high accuracy, creating privacy risks for users
behind nation-state, ISP, or enterprise monitoring.

A natural defense is to route traffic through an anonymizing proxy or middlebox that
\emph{re-shapes} metadata before it leaves the user's trust domain. Unlike
encryption, obfuscation does not hide that a flow exists; instead it degrades the
statistical fingerprints exploited by classifiers. The defender faces a multi-objective
problem: maximize privacy (reduce classifier confidence) while minimizing bandwidth
and latency overhead.

This paper presents an end-to-end measurement study on the CESNET-QUIC22 backbone
dataset~\cite{cesnet_quic22}. We make four contributions:
\begin{enumerate}
  \item A reproducible pipeline from raw QUIC flows to frozen deep learning attackers
        and deterministic obfuscators, with temporal train/test separation
        (weeks W-2022-44 / W-2022-45).
  \item A systematic evaluation of seven obfuscation policies with manifest-derived
        bandwidth and latency costs, plus formal Pareto dominance analysis.
  \item Evidence that \emph{jitter\_low} is a practical operating point: $\approx$1.7
        percentage point macro-F1 loss at 11\,ms mean latency overhead and zero
        bandwidth expansion---validated with bootstrap confidence intervals and paired
        significance tests.
  \item A comparison of CNN-BiLSTM vs.\ masked Transformer robustness under the same
        defenses, showing Transformers resist timing jitter better while both fail under
        aggressive MTU padding.
\end{enumerate}

Section~\ref{sec:related} reviews related work. Section~\ref{sec:threat} states the
threat model. Section~\ref{sec:methods} describes data, models, and obfuscators.
Section~\ref{sec:results} presents results; Section~\ref{sec:discussion} discusses
deployment implications; Section~\ref{sec:conclusion} concludes.

%==============================================================================
